\chapter{Abstract}
\begin{tabular}{@{} >{\bfseries}p{0.2\textwidth} p{0.77\textwidth} @{}}
\MakeUppercase{Keywords} &
TCP/IP; UDP sockets; reliability; checksum; sliding window; selective repeat;
adaptive timeout; congestion control; fast retransmit; fast recovery\\
\end{tabular}
\par

\begin{doublespacing}
This report describes how to design and build a TCP/IP-inspired communication
stack in C that works over UDP sockets. There are twenty-five sender-receiver
pairs in the project, from level 0 to level 24. It shows a clear progression
from simple bit transfer to more complex transport behavior that is similar to
real TCP principles. The first levels are all about encoding and decoding
messages. Middle levels add checking for parity, acknowledgments, retries based
on timeouts, and reliability through stop-and-wait. Later levels add things
like setting up and tearing down connections, packet checksums, sequence and
acknowledgment numbers, pipelining, sliding windows, chunk-based payload
transfer, out-of-order buffering, and adaptive retransmission timeout (RTO)
estimation.

The full project also includes the expanded web visualizer, which shows
versions 0 to 24. The code, finite state machine, and UML timeline are all
shown together in that interface. This makes it easier to learn about the
protocol progression as a whole rather than as separate source files.

A key goal of this project is to learn networking from the scratch through a
textbook-first approach. The conceptual basis follows to established
references, including the top-down networking approach by Kurose and Ross and
the implementation viewpoint of Stevens. AI tools were only used in a
responsible way to help with idea exploration, clarification, and debugging.
The final steps of implementation, coding, and validation were done by hand,
without copying any code.

\newpage

\begin{samepage}
The final levels implement congestion control ideas such as slow start,
congestion avoidance, and, in the most advanced version, fast retransmit and
fast recovery based on duplicate acknowledgements. This project is an
educational simplification of TCP and uses UDP as the underlying transport. It
captures many practical protocol engineering concepts, such as correctness,
efficiency, robustness to packet loss and reordering, and maintainable
protocol state management.
\end{samepage}

The educational value is very important. The versioned C code, textbook
references, and synchronized visualizer all work together to make TCP-like
behavior easier for new students to understand. These students often find
transport protocols and congestion control hard to understand at first.

The report talks about the architecture, algorithms, implementation choices,
version-by-version progress, and limitations. It also has diagrams, tables that
compare things, and analysis based on formulas for throughput and timeout
logic. The result is a full learning tool that connects networking theory with
real-world systems programming.
\end{doublespacing}
